\chapter{Polynomials}

\begin{center}\itshape
Among all monic polynomials of degree $n$, which one stays closest to $0$ on $[-1,1]$?\\
And how can the roots of a polynomial compute $1+\frac14+\frac19+\frac1{16}+\cdots$?
\end{center}

\section{Division, roots and multiplicities}

\begin{reminder}[polynomials over a field $\mathbb K$]
\begin{itemize}
\item \textbf{Euclidean division.} For $A,B\in\mathbb K[X]$ with $B\ne0$ there is a unique pair $(Q,R)$ with $A=BQ+R$ and $\deg R<\deg B$.
\item \textbf{Roots.} $a$ is a root of $P$ if and only if $(X-a)\mid P$. A non-zero polynomial of degree $n$ has at most $n$ roots, counted with multiplicity.
\item \textbf{Multiplicity} (in characteristic $0$). $a$ is a root of multiplicity at least $m$ if and only if $P(a)=P'(a)=\dots=P^{(m-1)}(a)=0$.
\item \textbf{D'Alembert--Gauss.} Every non-constant polynomial of $\C[X]$ has a root, so it splits into linear factors. Over $\R$, the irreducibles are the linear polynomials and the quadratics with negative discriminant.
\end{itemize}
\end{reminder}

\begin{proposition}[Vi\`ete's relations]
If $P=a_nX^n+a_{n-1}X^{n-1}+\dots+a_0=a_n\prod_{i=1}^n(X-x_i)$ with $a_n\ne0$, the elementary symmetric functions
$\sigma_k=\sum_{i_1<\dots<i_k}x_{i_1}\cdots x_{i_k}$ satisfy
\[
\sigma_k=(-1)^k\,\frac{a_{n-k}}{a_n}\qquad(1\le k\le n).
\]
\end{proposition}
\begin{proof}
Expand $\prod_i(X-x_i)$: the coefficient of $X^{n-k}$ is obtained by choosing $-x_i$ in exactly $k$ factors, so it equals $(-1)^k\sigma_k$.
\end{proof}

\section{Lagrange interpolation}

\begin{theorem}[Lagrange]\label{thm:lagrange}
Let $x_0,\dots,x_n\in\mathbb K$ be distinct and $y_0,\dots,y_n\in\mathbb K$. There is a unique $P\in\mathbb K[X]$ of degree at most $n$ with $P(x_i)=y_i$ for all $i$, namely
\[
P=\sum_{i=0}^ny_i\,L_i,\qquad L_i=\prod_{j\ne i}\frac{X-x_j}{x_i-x_j}.
\]
\end{theorem}
\begin{proof}
$L_i(x_j)$ is $1$ if $j=i$ and $0$ otherwise, so $P$ fits. If $\tilde P$ also fits, $P-\tilde P$ has degree at most $n$ and $n+1$ roots, hence is zero.
\end{proof}

\section{Chebyshev polynomials}

\begin{proposition}\label{prop:cheb}
For each $n\in\N$ there is a unique $T_n\in\R[X]$ with $T_n(\cos\theta)=\cos n\theta$ for all $\theta$. Moreover $T_0=1$, $T_1=X$,
\[
T_{n+1}=2X\,T_n-T_{n-1},
\]
$T_n$ has degree $n$ and leading coefficient $2^{n-1}$ (for $n\ge1$), and its roots are $\cos\frac{(2k+1)\pi}{2n}$, $0\le k\le n-1$.
\end{proposition}
\begin{proof}
From $\cos((n+1)\theta)+\cos((n-1)\theta)=2\cos\theta\cos n\theta$, the polynomials defined by the recurrence satisfy
$T_n(\cos\theta)=\cos n\theta$ by induction, and the degree and leading coefficient follow from the recurrence. Uniqueness: two such polynomials agree on $[-1,1]$, an infinite set.
The $n$ numbers $\cos\frac{(2k+1)\pi}{2n}$ are distinct and are roots since $\cos\frac{(2k+1)\pi}2=0$; as $\deg T_n=n$ there are no others.
\end{proof}

\section{Integer polynomials}

\begin{reminder}[irreducibility over $\Q$]
\begin{itemize}
\item \textbf{Rational roots.} If $p/q$ in lowest terms is a root of $a_nX^n+\dots+a_0\in\Z[X]$, then $p\mid a_0$ and $q\mid a_n$.
\item \textbf{Gauss' lemma.} A polynomial of $\Z[X]$ that factors in $\Q[X]$ factors in $\Z[X]$ with factors of the same degrees.
\item \textbf{Eisenstein's criterion.} Let $P=a_nX^n+\dots+a_0\in\Z[X]$ and $p$ a prime with $p\nmid a_n$, $p\mid a_i$ for $i<n$, and $p^2\nmid a_0$. Then $P$ is irreducible in $\Q[X]$.
\end{itemize}
\end{reminder}

\begin{proof}[Proof of Eisenstein]
By Gauss' lemma, suppose $P=BC$ in $\Z[X]$ with $\deg B,\deg C\ge1$. Reduce modulo $p$: $\bar P=\bar a_nX^n=\bar B\bar C$ in $(\Z/p\Z)[X]$.
Since $\Z/p\Z$ is a field, unique factorisation forces $\bar B=\beta X^k$, $\bar C=\gamma X^{n-k}$ with $k,n-k\ge1$ (degrees cannot drop because $p\nmid a_n$).
So $p$ divides the constant terms of both $B$ and $C$, and $p^2\mid a_0$: contradiction.
\end{proof}

\section{Exercises}

\begin{exercise}[$\star$]\label{ex:rem}
Find the remainder of the Euclidean division of $X^n$ by $(X-1)(X-2)$.
\end{exercise}

\begin{exercise}[$\star$]\label{ex:viete}
Let $x_1,x_2,x_3$ be the complex roots of $X^3-X+1$. Compute $x_1^2+x_2^2+x_3^2$ and $x_1^3+x_2^3+x_3^3$, without finding the roots.
\end{exercise}

\begin{exercise}[$\star\star$]\label{ex:simple}
Show that for $n\ge2$ the polynomial $X^n+X+1$ has only simple roots in $\C$.
\end{exercise}

\begin{exercise}[$\star\star$]\label{ex:lagsum}
Let $x_0,\dots,x_n$ be distinct complex numbers. Show that
\[
\sum_{i=0}^n\frac{x_i^k}{\prod_{j\ne i}(x_i-x_j)}=\begin{cases}0&\text{if }0\le k\le n-1,\\1&\text{if }k=n.\end{cases}
\]
\end{exercise}

\begin{exercise}[$\star\star$]\label{ex:eisen}
\begin{enumerate}[label=(\alph*)]
\item Show that $X^4+10X+5$ is irreducible over $\Q$.
\item Let $p$ be prime. Show that $\Phi_p=1+X+\dots+X^{p-1}$ is irreducible over $\Q$. \emph{Hint: look at $\Phi_p(X+1)$.}
\end{enumerate}
\end{exercise}

\begin{exercise}[$\star\star\star$]\label{ex:minimax}
\emph{Chebyshev's minimax property.} Show that every monic $P\in\R[X]$ of degree $n\ge1$ satisfies
$\sup_{[-1,1]}|P|\ge2^{1-n}$, with equality for $P=2^{1-n}T_n$.
\end{exercise}

\begin{exercise}[$\star\star\star$]\label{ex:basel}
\emph{The Basel problem through roots.} Let $m\ge1$ and $N=2m+1$.
\begin{enumerate}[label=(\alph*)]
\item Using $\sin N\theta=\operatorname{Im}\left((\cos\theta+i\sin\theta)^N\right)$, show that for $0<\theta<\frac\pi2$,
\[
\sin N\theta=\sin^N\theta\sum_{j=0}^m(-1)^j\binom N{2j+1}\cot^{2(m-j)}\theta.
\]
\item Deduce that the numbers $\cot^2\frac{k\pi}N$, $1\le k\le m$, are the roots of $\sum_{j=0}^m(-1)^j\binom N{2j+1}X^{m-j}$, and that
$\displaystyle\sum_{k=1}^m\cot^2\frac{k\pi}N=\frac{m(2m-1)}3$.
\item Using $\cot^2x<\frac1{x^2}<1+\cot^2x$ on $\left]0,\frac\pi2\right[$, prove that $\displaystyle\sum_{k=1}^\infty\frac1{k^2}=\frac{\pi^2}6$.
\end{enumerate}
\end{exercise}

\begin{exercise}[$\star\star\star$]\label{ex:sos}
Let $P\in\R[X]$ with $P(x)\ge0$ for all $x\in\R$. Show that there are $A,B\in\R[X]$ with $P=A^2+B^2$.
\end{exercise}

\begin{beauty}
Exercise~\ref{ex:basel} computes an infinite sum with nothing but Vi\`ete's relations; Tome~4 gives five more proofs.
Exercise~\ref{ex:eisen}(b) is the first step towards the constructible polygons of Tome~2, where the irreducibility of
$\Phi_{17}$ explains Gauss' 17-gon. Chebyshev's minimax property is why numerical analysts interpolate at the roots of
$T_n$ rather than at equally spaced points, as the notebook shows.
\end{beauty}

\section{Solutions}

\begin{solution}{ex:rem}
Write $X^n=(X-1)(X-2)Q+aX+b$. Evaluating at $1$ and $2$: $a+b=1$ and $2a+b=2^n$. So the remainder is $(2^n-1)X+2-2^n$.
\end{solution}

\begin{solution}{ex:viete}
Here $\sigma_1=0$, $\sigma_2=-1$, $\sigma_3=-1$. Then $\sum x_i^2=\sigma_1^2-2\sigma_2=2$. Each root satisfies $x_i^3=x_i-1$, so
$\sum x_i^3=\sum x_i-3=-3$.
\end{solution}

\begin{solution}{ex:simple}
A multiple root $a$ satisfies $a^n+a+1=0$ and $na^{n-1}+1=0$. Multiplying the second by $a$, $na^n=-a$, and substituting $a^n=-a-1$ gives
$-n(a+1)=-a$, so $a=-\frac n{n-1}$. Then $|a|=\frac n{n-1}>1$, so $|a|^{n-1}>1$; but $na^{n-1}=-1$ forces $|a|^{n-1}=\frac1n<1$. Contradiction.
\end{solution}

\begin{solution}{ex:lagsum}
For $k\le n$, the polynomial $X^k$ has degree at most $n$, so by uniqueness in Theorem~\ref{thm:lagrange}, $X^k=\sum_ix_i^kL_i$.
The coefficient of $X^n$ in $L_i$ is $1/\prod_{j\ne i}(x_i-x_j)$, so the coefficient of $X^n$ on the right is exactly the sum in question.
On the left it is $0$ for $k<n$ and $1$ for $k=n$.
\end{solution}

\begin{solution}{ex:eisen}
(a) Eisenstein with $p=5$: $5\nmid1$, $5$ divides $0,0,10,5$, and $25\nmid5$.

(b) $\Phi_p=\frac{X^p-1}{X-1}$, so $\Phi_p(X+1)=\frac{(X+1)^p-1}X=\sum_{k=1}^p\binom pkX^{k-1}$. It is monic, its other coefficients
$\binom pk$ ($1\le k\le p-1$) are divisible by $p$ (Exercise~\ref{ex:fermat}), and its constant term is $\binom p1=p$, not divisible by $p^2$.
So $\Phi_p(X+1)$ is irreducible by Eisenstein, and so is $\Phi_p$, since a factorisation $\Phi_p=BC$ would give $\Phi_p(X+1)=B(X+1)C(X+1)$.
\end{solution}

\begin{solution}{ex:minimax}
For $P=2^{1-n}T_n$: it is monic by Proposition~\ref{prop:cheb}, and $|T_n(\cos\theta)|=|\cos n\theta|\le1$ with equality at $\theta=0$, so $\sup|P|=2^{1-n}$.

Now suppose $P$ is monic of degree $n$ with $\sup_{[-1,1]}|P|<2^{1-n}$, and let $Q=2^{1-n}T_n-P$, of degree at most $n-1$ (the leading terms cancel).
At $y_k=\cos\frac{k\pi}n$, $0\le k\le n$, we have $2^{1-n}T_n(y_k)=2^{1-n}(-1)^k$, so $Q(y_k)$ has the sign of $(-1)^k$, strictly.
By the intermediate value theorem $Q$ vanishes in each of the $n$ intervals $]y_{k+1},y_k[$, so $Q$ has $n$ roots and degree at most $n-1$: $Q=0$,
which contradicts $\sup|P|<2^{1-n}=\sup|2^{1-n}T_n|$.
\end{solution}

\begin{solution}{ex:basel}
(a) By the binomial theorem, the imaginary part of $(\cos\theta+i\sin\theta)^N$ collects the odd powers of $i\sin\theta$:
\[
\sin N\theta=\sum_{j=0}^m\binom N{2j+1}(-1)^j\cos^{N-2j-1}\theta\,\sin^{2j+1}\theta
=\sin^N\theta\sum_{j=0}^m(-1)^j\binom N{2j+1}\cot^{N-2j-1}\theta,
\]
and $N-2j-1=2(m-j)$.

(b) For $\theta_k=\frac{k\pi}N$ with $1\le k\le m$ we have $0<\theta_k<\frac\pi2$, $\sin N\theta_k=0$ and $\sin\theta_k\ne0$, so $\cot^2\theta_k$ is a root of
$R=\sum_j(-1)^j\binom N{2j+1}X^{m-j}$. These $m$ numbers are distinct since $\cot^2$ is injective on $\left]0,\frac\pi2\right[$, and $\deg R=m$, so they are all its roots.
By Vi\`ete,
\[
\sum_{k=1}^m\cot^2\theta_k=\frac{\binom N3}{\binom N1}=\frac{(2m+1)(2m)(2m-1)}{6(2m+1)}=\frac{m(2m-1)}3.
\]
(c) The inequality follows from $0<\sin x<x<\tan x$ on $\left]0,\frac\pi2\right[$, since $\frac1{\sin^2x}=1+\cot^2x$.
Apply it at $x=\theta_k$ and sum over $1\le k\le m$:
\[
\frac{m(2m-1)}3<\frac{N^2}{\pi^2}\sum_{k=1}^m\frac1{k^2}<m+\frac{m(2m-1)}3.
\]
Multiply by $\frac{\pi^2}{N^2}=\frac{\pi^2}{(2m+1)^2}$. Both bounds tend to $\frac{\pi^2}6$ as $m\to\infty$ (they are equivalent to $\frac{\pi^2}{4m^2}\cdot\frac{2m^2}3$), so $\sum\frac1{k^2}=\frac{\pi^2}6$.
\end{solution}

\begin{solution}{ex:sos}
If $P=0$ take $A=B=0$. Otherwise, a real root of odd multiplicity would make $P$ change sign there, so every real root has even multiplicity;
and the leading coefficient $c$ is positive since $P(x)\to\pm\infty$ with the sign of $c$ (the degree is even). Over $\C$, the non-real roots come in conjugate pairs, so
\[
P=c\,R^2\prod_{j=1}^s(X-z_j)(X-\bar z_j),\qquad R=\prod_i(X-a_i)^{m_i}\in\R[X].
\]
Write $\prod_j(X-z_j)=U+iV$ with $U,V\in\R[X]$. Then $\prod_j(X-\bar z_j)=U-iV$ and the product is $U^2+V^2$.
So $P=(\sqrt c\,RU)^2+(\sqrt c\,RV)^2$.
\end{solution}
